\section{Conclusión}\label{sec:conclusion_informe_final}

La motivación de este proyecto nació de la observación directa de las ineficiencias en los procesos de gestión del LabDIC, donde la experiencia sostenida en el laboratorio permitió identificar con claridad los puntos críticos que una solución tecnológica podía resolver. La combinación entre formación académica e investigación aplicada fundamentó la selección de un stack tecnológico que satisface los requerimientos actuales y sienta las bases para un eventual despliegue en producción.

El modelo de base de datos evolucionó desde una concepción inicial de tres ejes funcionales hacia una arquitectura de seis módulos cohesivos, donde la modularidad adoptada es de naturaleza conceptual: no todos los módulos requieren tablas propias, sino que algunos se expresan a través de campos y relaciones dentro de tablas existentes. La integración de SQLAlchemy y Alembic proporcionó un marco robusto para la evolución iterativa del esquema, asegurando que cada cambio quedara documentado y fuera reproducible.

La documentación arquitectónica del sistema evidenció que las decisiones de diseño responden a principios consistentes de separación de responsabilidades aplicados de forma simétrica en ambas capas. En el backend, la división entre DTOs, repositorios y controladores garantiza que cada componente tenga una responsabilidad acotada; en el frontend, la cadena \texttt{types/} $\rightarrow$ \texttt{services/} $\rightarrow$ stores de Pinia replica esta misma lógica. Esta uniformidad hace que añadir un nuevo módulo implique replicar un patrón ya establecido.

La implementación del backend consolidó la API REST como una capa completa y estructurada. Los módulos más complejos, Dispositivos y Préstamos, ilustraron cómo la lógica de negocio puede encapsularse en el repositorio dejando los controladores livianos. Las migraciones generadas durante este proceso confirman que el desarrollo del backend no es un evento único, sino un proceso iterativo que evoluciona en paralelo con las demás capas del sistema.

El frontend completó el ciclo de desarrollo mediante seis fases progresivas que construyeron la interfaz de forma incremental. La interdependencia real entre frontend y backend quedó en evidencia cuando la implementación de vistas reveló la necesidad de correcciones en cascada en la base de datos. Desde la perspectiva del usuario, el sistema resultante ofrece una experiencia diferenciada según el rol, con soporte para modo claro y oscuro, diseño \textit{responsive} y retroalimentación visual mediante notificaciones. Con todas las capas implementadas, la documentación del trabajo futuro identifica mejoras pendientes y propone nuevos módulos que la arquitectura modular de la plataforma está preparada para incorporar.